\slajdid{09_1-s037}
\begin{frame}[label=09_1-s037]
\gusttekst\setlength{\abovedisplayskip}{4pt}\setlength{\belowdisplayskip}{4pt}
\begin{columns}[c,onlytextwidth]\column{.23\textwidth}\centering\input{slike/tikz/s037_invertor_sa_aktivnim_opterecenjem.tex}\column{.75\textwidth}
Da pogledamo režim rada tranzistora $T_L$. Naponi na tranzistoru su
\[V_{GS,L}=V_{G,L}-V_{S,L}=V_C-V_O\]
\[V_{DS,L}=V_{D,L}-V_{S,L}=V_{DD}-V_O\]
Uslov da radi u zasićenju
\[V_{DS,L}\geq V_{GS,L}-V_{Tn,L}\]
\[V_{DD}-V_O\geq V_C-V_O-V_{Tn,L}\]
\[V_{DD}+V_{Tn,L}\geq V_C\]
i dobija se zanimljiv rezultat. Vidi se da se podešavanjem kontrolnog napona $V_C$ može postići da tranzistor $T_L$ uvek radi u zasićenju
\[V_C<V_{DD}+V_{Tn,L}\]
odnosno da uvek radi u omskoj oblasti
\[V_C>V_{DD}+V_{Tn,L}\]
\end{columns}
\end{frame}
